\documentclass[10pt]{article}
\usepackage[a4paper,margin=25mm]{geometry}
\usepackage{amsmath,amssymb,amsthm,booktabs,array}
\usepackage[T1]{fontenc}
\usepackage{lmodern}
\usepackage[colorlinks=true,urlcolor=blue,citecolor=blue,linkcolor=blue]{hyperref}
\newtheorem{theorem}{Theorem}
\newtheorem{lemma}{Lemma}
\newtheorem{corollary}{Corollary}
\newcommand{\F}{\mathbb F}
\newcommand{\C}{\mathbb C}
\title{Hamming extensions and a short proof from Min\'a\v{c} et al.'s cycle\newline
answering Nguyen and T\^an's primeness question}
\author{Haobo Ma\\{\small ChronoAI Pte Ltd, Singapore}
\and Wenlin Zhang\\{\small National University of Singapore}\\
{\small ChronoAI Pte Ltd, Singapore}}
\date{November 19, 2026}
\begin{document}
\maketitle
\begin{abstract}
Using the cycle example of Min\'a\v{c}, Muller, Nguyen and T\^an
(2022, Corollary 2.5), we give a short negative answer to Nguyen and
T\^an's Question 4.17 (arXiv:2508.10348v1, Section 4.3, p.~19):
$\Gamma(\mathbb Z/8\mathbb Z,\{1,-1\})=C_8$ is prime, connected and
anti-connected, with adjacency nullity two. More generally, $C_n$ has
these graph properties for $n\geq5$ and is singular exactly when $4\mid n$.
Extending Chen's binary polynomial construction with the tensor spectrum
described by Sander, we realize every Hamming graph $H(d,q)$ over a
prime-power alphabet as a subgroup-unit Cayley graph of a Frobenius group
algebra. For $d\geq2$, except $(d,q)=(2,2)$, it is prime and
anti-connected; it is singular exactly when $q\mid d$, with nullity
$\binom d{d(q-1)/q}(q-1)^{d(q-1)/q}$.
These examples also answer Question 2.20 of arXiv:2603.21239v1 negatively.
Finally, a short proof of Hendrickson's duality answers the
supercharacter-axiom question in Remark 2.2 of the first source affirmatively.
\end{abstract}

\section{The question and an eight-vertex answer}
For a finite commutative ring $R$ with identity and a subgroup
$U\leq R^\times$ containing $-1$, let $\Gamma(R,U)$ be the simple
graph on $R$ in which distinct $a,b$ are adjacent when $a-b\in U$.
A graph is \emph{anti-connected} if its complement is connected.
A vertex set $M$ is \emph{homogeneous} if every vertex outside $M$
is adjacent to all or none of $M$. A graph is \emph{prime} if no
homogeneous set has size between $2$ and $|V|-1$.
A finite ring of characteristic $n$ is \emph{Frobenius} in the convention
of~\cite[Section 4.1]{NT} if it admits a $\mathbb Z/n\mathbb Z$-linear
functional to $\mathbb Z/n\mathbb Z$ whose kernel contains no nonzero ideal.

For Frobenius rings, Nguyen and T\^an~\cite[Question 4.17]{NT} ask,
assuming that the graph and its complement are connected, whether it is prime
if and only if zero is absent from its adjacency spectrum. Their later
Question 2.20~\cite[Section 2.2.4, p.~9]{NT26} asks whether primeness
implies nonsingularity for finite rings. Both answers are no, already
for the eight-vertex cycle. This example is the generalized Paley graph
$P_8$ in Min\'a\v{c} et al.~\cite[Corollary 2.5 and Theorem 3.5]{MMNT};
those results identify the cycle and its spectrum. The primeness
argument below is self-contained. Alternatively, their nullity two,
together with Chudnovsky et al.'s nullity bound~\cite[Corollary 4.3]{CC},
gives the same conclusion. Thus the cycle answer follows from earlier
constructions and results, rather than requiring a new graph.

\begin{lemma}\label{lem:twins}
A connected triangle-free graph with no distinct vertices having the
same open neighborhood is prime.
\end{lemma}
\begin{proof}
If a proper homogeneous set $M$ contains an edge, every outside vertex
adjacent to any member of $M$ is adjacent to both ends of that edge,
creating a triangle. Hence there are no edges between $M$ and its
complement, contrary to connectedness. Thus a proper homogeneous set
is independent. If it has two members, they have the same outside
neighbors and no inside neighbors, contradicting the hypothesis.
\end{proof}

\begin{theorem}\label{thm:cycles}
For every $n\geq5$, $R=\mathbb Z/n\mathbb Z$ is Frobenius and
$U=\{1,-1\}$ is a subgroup of its units. The graph $\Gamma(R,U)=C_n$
is connected, anti-connected and prime. Its adjacency nullity is two
when $4\mid n$, and zero otherwise. In particular, every $C_{4m}$,
$m\geq2$, answers both primeness questions negatively.
\end{theorem}
\begin{proof}
The identity functional on $R$ has zero kernel, proving the Frobenius
property. The graph is the cycle with edges $a\sim a\pm1$.
In its complement, each original edge $a\sim a+1$ can be replaced by
the path $a,a+3,a+1$, since neither $2$ nor $3$ is $\pm1$ modulo
$n\geq5$. Thus the complement is connected.

The cycle is triangle-free and has no equal open neighborhoods for
$n\geq5$. Indeed, if a vertex $b\neq0$ had neighborhood
$\{-1,1\}$, then $b=2$ or $b=-2$ because it is adjacent to $1$ or
$-1$. Equality of both neighbors would force $n=4$.
Lemma~\ref{lem:twins} therefore proves primeness.

For a complex vector $v=(v_a)_{a\in\mathbb Z/n\mathbb Z}$, the
equation $Av=0$ is $v_{a+1}=-v_{a-1}$. All solutions on the integer
line have the repeating pattern $(s,t,-s,-t)$.
If $4\mid n$, both parameters are free. If $n\equiv2\pmod4$,
periodicity forces $s=-s$ and $t=-t$; if $n$ is odd, it forces both
$s=t$ and $s=-t$. Over $\C$ these conditions give $s=t=0$.
This proves the nullity statement. For $n=8$, the two repeating
patterns $(1,0,-1,0)$ and $(0,1,0,-1)$ give an explicit nullspace
basis and rank $8-2=6$.
\end{proof}

\section{A Hamming family over every finite field}
The cube construction also fits a larger family. Let $q=p^r$ be a
prime power, let $\F_q$ be the field with $q$ elements, and let $G$
be a finite abelian group of order $d$, written additively with identity
$0_G$. Define the group algebra
\[
 R=\F_q[G]=\left\{\sum_{g\in G}a_g e_g:a_g\in\F_q\right\},
 \qquad e_g e_h=e_{g+h},\qquad 1_R=e_{0_G},
\]
with coefficientwise addition and bilinear multiplication. Put
\[
 U=\{c e_g:c\in\F_q^\times,\ g\in G\}.
\]
The inverse of $ce_g$ is $c^{-1}e_{-g}$; hence $U$ is a subgroup of
$R^\times$ of order $d(q-1)$ containing $-1_R$.
The \emph{Hamming graph} $H(d,q)$ has vertex set $\F_q^G$ and
edges joining vectors differing in exactly one coordinate. Coefficients
identify $\Gamma(R,U)$ with $H(d,q)$.

For $q=2$ and cyclic $G$, this is Chen's polynomial construction
$\F_2[x]/(x^d-1)$ with connection set $\{1,x,\ldots,x^{d-1}\}$:
take the ideal $(x^d+1)$ in~\cite[Section 2, equation (4), p.~5]{Chen}.
The coefficient map makes its graph the cube $Q_d=H(d,2)$.
The tensor spectral method below is the one described in
Sander~\cite[Theorem 2.2 and Lemma 3.1]{Sander}; its cubelike
character version appears in Cheung--Godsil~\cite[Section 3]{CG}.
Nguyen's recent scalar-orbit character formula~\cite[Proposition 2.2]{Matrix}
also applies to this scalar-stable connection set.
The extension here is the scalar-basis unit subgroup over every
$\F_q$, with an exact singularity criterion and its application
to the two questions.

\begin{theorem}\label{thm:hamming}
Let $q$ be a prime power and $G$ a finite abelian group of order $d$.
The ring $R=\F_q[G]$ is Frobenius, and $\Gamma(R,U)\cong H(d,q)$
has adjacency eigenvalues
\[
 \lambda_k=d(q-1)-qk,\qquad
 \text{multiplicity }b_k=\binom dk(q-1)^k,\qquad 0\leq k\leq d.
\]
If $d\geq2$ and $(d,q)\neq(2,2)$, this graph is connected,
anti-connected and prime. It is singular exactly when $q\mid d$,
and then its nullity is $b_{d(q-1)/q}<q^d/2$.
\end{theorem}
\begin{proof}
Choose a nonzero $\F_p$-linear map $\ell:\F_q\to\F_p$ and set
$\psi(a)=\ell(a_{0_G})$. Given $a\neq0$, choose $g$ with
$a_g\neq0$, and choose $c\in\F_q$ such that $\ell(ca_g)\neq0$;
this is possible since multiplication by $a_g$ is bijective.
Then $\psi(ce_{-g}a)\neq0$. Thus no nonzero ideal lies in
$\ker\psi$, as required.

Write $K_q$ for the complete graph on $\F_q$, with adjacency matrix
$B$. Its constant functions have eigenvalue $q-1$, and its
$(q-1)$-dimensional subspace of functions with sum zero has
eigenvalue $-1$, since
$(Bf)(a)=\sum_b f(b)-f(a)$. The adjacency matrix of $H(d,q)$
is the sum of $d$ tensor operators, each applying $B$ in one
coordinate and the identity in the others. Choose, in each coordinate,
the constant function and a basis of the zero-sum subspace. Their
tensor products form a basis on $\F_q^G$. A product using $k$
zero-sum factors has eigenvalue $(d-k)(q-1)-k$, and there are
$\binom dk(q-1)^k$ such products. This proves the entire spectrum.
Zero occurs exactly when $k=d(q-1)/q$ is an integer, equivalently
$q\mid d$, because $\gcd(q,q-1)=1$.

Changing one coordinate at a time proves connectedness. For $d\geq3$,
a coordinate edge from $v$ to $w$ can be replaced in the complement
by $v,z,w$, where $z$ differs from $v$ in two coordinates other
than the changed coordinate. The two distances are two and three.
For $d=2$ and $q\geq3$, take $z$ different from both endpoint
values in the changed coordinate and different from their common
value in the other coordinate. Both distances are two. Hence the
complement is connected in the stated range.

To prove primeness, use the nullity bound of Chudnovsky et al.
\cite[Corollary 4.3]{CC}: a nonprime connected, anti-connected
Cayley graph on a commutative ring with its connection subgroup
in the units has nullity at least half its order. If the graph is
nonsingular, that bound already proves primeness. Otherwise write
$d=qm$ and $k=(q-1)m$. The adjacent multiplicities satisfy
\[
 \frac{b_{k-1}}{b_k}=\frac{m}{m+1},\qquad
 \frac{b_{k+1}}{b_k}=\frac{(q-1)m}{(q-1)m+1}.
\]
Their sum exceeds one unless $(q,m)=(2,1)$, precisely the excluded
case $(d,q)=(2,2)$. Consequently
$q^d=\sum_j b_j\geq b_{k-1}+b_k+b_{k+1}>2b_k$.
The bound therefore proves primeness also in the singular case.
\end{proof}

The primeness of these Hamming graphs also follows from the later
Cartesian-product theorem of El Amri, Boudabbous and
Yaich~\cite[Theorem 1.2]{Cartesian}: a Cartesian product of two
nontrivial connected graphs is prime if one factor has at least three
vertices. Apply it to $K_q\mathbin\square H(d-1,q)$; for $q=2$
and $d\geq3$, the second factor has at least four vertices.

\begin{corollary}\label{cor:examples}
For every prime power $q$, the family in Theorem~\ref{thm:hamming}
gives prime singular examples whenever $q\mid d$ and $(d,q)\neq(2,2)$.
In particular, the binary examples are $Q_d$ for even $d\geq4$;
the ternary example $H(3,3)$ has $27$ vertices, rank $15$ and
nullity $12$. Odd prime powers give examples in odd characteristic
and of odd vertex order.
\end{corollary}
\begin{proof}
Apply the theorem. For $q=d=3$, the spectrum is
$6,3,0,-3$ with multiplicities $1,6,12,8$.
For every odd $q$, take $d=q$, so $q^d$ is odd.
\end{proof}

For cyclic $G$, the map $x\mapsto e_1$ identifies the ring with
$\F_q[x]/(x^d-1)$: the $d$ residue monomials map to the $d$
group-algebra basis elements. Thus the binary case $d=4$ is
$R=\F_2[x]/(x^4-1)$, $U=\{1,x,x^2,x^3\}$, with spectrum
$4,2,0,-2,-4$ of multiplicities $1,4,6,4,1$ and nullity six.
Here $Q_4$ is the first qualifying even cube, whereas $C_8$ above
has fewer vertices. No minimality among all ring graphs is needed.
For the boundary cases, $H(1,q)$ is complete, and the complement
of $H(2,2)=C_4$ is disconnected; its opposite pairs are homogeneous.

The hypotheses of the source's positive cases remain important.
For example, the singular examples have $p\mid d(q-1)=|U|$,
so they do not satisfy~\cite[Proposition 4.20]{NT}.
For a converse not asserted there, divisibility $p\mid|U|$ does
not imply singularity: $q=4,d=2$ gives $|U|=6$, characteristic two,
and spectrum $6,2,-2$, with multiplicities $1,6,9$.

\section{Hendrickson duality and the neighboring axiom question}
The question in~\cite[Remark 2.2, p.~4]{NT} asks whether condition
(4) of its Definition 2.1 follows from the other conditions. It does;
this is the duality in Hendrickson~\cite[Theorem 10.3 and the paragraph
following Corollary 10.5]{Hendrickson}. We include the short Fourier proof.

\begin{corollary}\label{cor:duality}
Let $A$ be a finite abelian group, and let $\widehat A$ be its
complex character group. Suppose $\mathcal K$ partitions $A$,
$\mathcal X$ partitions $\widehat A$, $|\mathcal K|=|\mathcal X|$,
and $\sigma_X=\sum_{\chi\in X}\chi$ is constant on each
$K\in\mathcal K$ for every $X\in\mathcal X$.
Then, for each $K\in\mathcal K$, the sum $\sum_{a\in K}\chi(a)$
is independent of $\chi\in X$ for every fixed $X\in\mathcal X$.
In particular, Remark 2.2 is answered affirmatively, even without
the additional symmetry assumption on $\mathcal K$.
\end{corollary}
\begin{proof}
Let $V$ be the space of complex functions constant on every block
of $\mathcal K$. It has dimension $|\mathcal K|$.
The functions $\sigma_X$ lie in $V$ and are linearly independent:
in the character basis their nonempty supports are disjoint.
They therefore form a basis of $V$. For the indicator $1_K$,
write $1_K=\sum_X c_X\sigma_X$.
With inner product
$\langle f,h\rangle=|A|^{-1}\sum_{a\in A}f(a)\overline{h(a)}$,
character orthogonality gives, for $\chi\in X$,
\[
 \frac1{|A|}\sum_{a\in K}\overline{\chi(a)}
 =\langle1_K,\chi\rangle=c_X.
\]
Conjugating proves the required constancy. Neither singleton-zero
nor symmetry of the blocks is used in this argument.
\end{proof}

\section{Coverage and mathematical evidence}
\noindent Unqualified source items in the following table refer to~\cite{NT}.
\begin{center}\small
\begin{tabular}{@{}p{.25\linewidth}p{.69\linewidth}@{}}
\toprule
Source item & Coverage \\
\midrule
Question 4.17 & Answered: no. Theorem~\ref{thm:cycles} gives $C_8$;
Theorem~\ref{thm:hamming} gives the prime-power family. The cycle and
binary constructions have prior sources~\cite{MMNT,Chen}.\\
\cite{NT26}, Question 2.20 & Answered: no, by the same commutative examples.\\
Remark 2.2 & Answered: yes. Reproved in Corollary~\ref{cor:duality};
prior proof: Hendrickson (2009), Theorem 10.3.\\
Proposition 4.16; \cite{NT26}, Proposition 2.19 & Prior proof cited;
the nonprime-to-singular implication is unchanged.\\
Proposition 4.18; \cite{NT26}, Proposition 2.21 & Prior proof cited;
the full-unit-group case is unchanged.\\
Proposition 4.20 & Prior proof cited. Its coprime-order hypothesis
excludes our singular Hamming graphs. A converse is not asserted in
the source; it fails for $H(2,4)$.\\
Propositions 4.21--4.22 & Not covered beyond their cited specialized
positive cases; their Jacobi-symbol and power-unit hypotheses are
not replaced by arbitrary subgroups.\\
Prime-power Hamming family & Not posed; proved in
Theorem~\ref{thm:hamming}, using the prior tensor spectrum and
nullity bound~\cite{Sander,CC}.\\
\bottomrule
\end{tabular}
\end{center}

\subsection*{Written, formal and finite scope}
The cycle proof, prime-power Hamming theorem, complete spectra and
multiplicities, polynomial descriptions and duality corollary have the
written proofs above. The
\href{https://github.com/the-omega-institute/trureturing/blob/2eae85555e169007ff5548c46eb63cd2408f877f/D5/S3/ArithUnits/NguyenTanUnitaryCayleyPrimeSpectrum.lean#L669}{pinned Lean declaration
\texttt{result}} states the negation of the affirmative form of
Question 4.17. Its
\href{https://github.com/the-omega-institute/trureturing/blob/2eae85555e169007ff5548c46eb63cd2408f877f/D5/S3/ArithUnits/NguyenTanUnitaryCayleyPrimeSpectrum.lean#L631}{supporting lemma
\texttt{group\_algebra\_counterexample}} covers the binary Frobenius
group algebra, both connectivity properties, primeness and a zero
eigenvalue for finite abelian groups of even order at least four;
\href{https://github.com/the-omega-institute/trureturing/blob/2eae85555e169007ff5548c46eb63cd2408f877f/D5/S3/ArithUnits/NguyenTanUnitaryCayleyPrimeSpectrum.lean#L650}{\texttt{even\_cyclic\_family}}
specializes to cyclic groups. The exported contradiction uses order four.
The formal statement is restricted to this binary family; it does not
include the cycle construction, the prime-power extension, complete
multiplicities or the duality corollary. The accompanying
\texttt{verify.py} gives exact finite checks of cycles, tensor
eigenvectors, explicit rings and character partitions, with parameter
ranges and results in \texttt{verification.json}; these finite checks
are separate from the general proofs.

\begin{thebibliography}{10}
\small
\setlength{\itemsep}{2pt}
\bibitem{NT} T.~T.~Nguyen and N.~D.~T\^an,
\emph{Supercharacters of finite abelian groups and applications to spectra
of $U$-unitary Cayley graphs},
\href{https://arxiv.org/abs/2508.10348v1}{arXiv:2508.10348v1} (2025).
\bibitem{NT26} T.~T.~Nguyen and N.~D.~T\^an,
\emph{On $U$-unitary Cayley graphs over finite rings},
\href{https://arxiv.org/abs/2603.21239v1}{arXiv:2603.21239v1} (2026).
\bibitem{MMNT} J.~Min\'a\v{c}, L.~Muller, T.~T.~Nguyen and N.~D.~T\^an,
\emph{On the Paley graph of a quadratic character},
\href{https://arxiv.org/abs/2212.02005v1}{arXiv:2212.02005v1} (2022),
Corollary 2.5 and Theorem 3.5; the cycle statement is Corollary 2.6 in
\href{https://arxiv.org/abs/2212.02005v2}{version 2} (2023).
\bibitem{Chen} H.~Chen,
\emph{Quotients of polynomial rings and regular $t$-balanced Cayley maps
on abelian groups},
\href{https://arxiv.org/abs/1407.5295v5}{arXiv:1407.5295v5} (2017),
Section 2, equation (4); first version (2014), Notation 2.2.
\bibitem{Sander} T.~Sander,
\emph{Eigenspaces of Hamming graphs and unitary Cayley graphs},
Ars Math. Contemp. \textbf{3} (2010), 13--19,
\href{https://doi.org/10.26493/1855-3974.100.7f8}{doi:10.26493/1855-3974.100.7f8}.
\bibitem{CG} W.-C.~Cheung and C.~Godsil,
\emph{Perfect state transfer in cubelike graphs},
Linear Algebra Appl. \textbf{435} (2011), 2468--2474,
\href{https://doi.org/10.1016/j.laa.2011.04.022}{doi:10.1016/j.laa.2011.04.022};
\href{https://arxiv.org/abs/1010.4721v2}{arXiv:1010.4721v2}, Section 3.
\bibitem{CC} M.~Chudnovsky, M.~Cizek, L.~Crew, J.~Min\'a\v{c},
T.~T.~Nguyen, S.~Spirkl and N.~D.~T\^an,
\emph{On prime Cayley graphs},
\href{https://arxiv.org/abs/2401.06062v1}{arXiv:2401.06062v1} (2024),
Corollary 4.3.
\bibitem{Hendrickson} A.~O.~F.~Hendrickson,
\emph{Construction of supercharacter theories of finite groups},
\href{https://arxiv.org/abs/0905.3538v1}{arXiv:0905.3538v1} (2009),
Theorem 10.3, p.~19, and the paragraph following Corollary 10.5, p.~20.
\bibitem{Cartesian} M.~El Amri, Y.~Boudabbous and N.~Yaich,
\emph{On primality of Cartesian product of graphs},
Arab J. Math. Sci. \textbf{31} (2025), 143--152
(published online July 16, 2024),
\href{https://doi.org/10.1108/AJMS-09-2023-0023}{doi:10.1108/AJMS-09-2023-0023}.
\bibitem{Matrix} T.~T.~Nguyen,
\emph{Spectra of gcd-graphs over a matrix ring},
\href{https://arxiv.org/abs/2611.15026v2}{arXiv:2611.15026v2} (2026),
Proposition 2.2.
\end{thebibliography}
\medskip\noindent\small
\copyright\ 2026 Haobo Ma and Wenlin Zhang.
Licensed under \href{https://creativecommons.org/licenses/by/4.0/}{CC BY 4.0}.
\end{document}
